%!TEX program = xelatex
%!TEX root = ../all_beamer.tex

%\documentclass[10pt,aspectratio=169]{beamer}
%
%% ── Theme ────────────────────────────────────────────────────────────────────
%\usetheme{Madrid}
%\usecolortheme{default}
%
%% ── Custom Tor Vergata palette ───────────────────────────────────────────────
%\definecolor{MainBlue}{HTML}{1F3C6E}
%\definecolor{AccentBlue}{HTML}{0070C0}
%\definecolor{torgray}{HTML}{E8EDF4}
%\definecolor{AlertRed}{HTML}{C0392B}
%\definecolor{ExGreen}{HTML}{1A7A4A}
%\definecolor{warnOrange}{HTML}{D35400}
%
%\setbeamercolor{palette primary}{bg=MainBlue,fg=white}
%\setbeamercolor{palette secondary}{bg=MainBlue!80,fg=white}
%\setbeamercolor{palette tertiary}{bg=MainBlue!60,fg=white}
%\setbeamercolor{palette quaternary}{bg=MainBlue,fg=white}
%\setbeamercolor{structure}{fg=MainBlue}
%\setbeamercolor{section in toc}{fg=MainBlue}
%\setbeamercolor{block title}{bg=MainBlue,fg=white}
%\setbeamercolor{block body}{bg=torgray}
%\setbeamercolor{block title alerted}{bg=AlertRed,fg=white}
%\setbeamercolor{block body alerted}{bg=AlertRed!10}
%\setbeamercolor{block title example}{bg=ExGreen,fg=white}
%\setbeamercolor{block body example}{bg=ExGreen!10}
%\setbeamercolor{alerted text}{fg=AlertRed}

% ── Packages ─────────────────────────────────────────────────────────────────
%\usepackage{amsmath,amssymb,bm}
%\usepackage{pgfplots}
%\pgfplotsset{compat=1.18}
%\usepackage{tikz}
%\usetikzlibrary{arrows.meta,positioning,shapes,decorations.pathreplacing,calc,decorations.markings}
%\usepgfplotslibrary{fillbetween}
%\usepackage{booktabs}
%\usepackage{multicol}
%\usepackage{graphicx}
%\usepackage{xcolor}
%\usepackage{mathtools}
%
%% ── Math macros ──────────────────────────────────────────────────────────────
%
%
%\newcommand{\ow}{\overline{\mathbf{w}}}
%\newcommand{\ox}{\overline{\mathbf{x}}}
%\newcommand{\oX}{\overline{\mathbf{X}}}
%
%\newcommand{\argmin}[1]{\underset{#1}{\operatorname{arg\,min}}}
%\newcommand{\argmax}[1]{\underset{#1}{\operatorname{arg\,max}}}
%\DeclareMathOperator{\Tr}{Tr}

% ── Footline ─────────────────────────────────────────────────────────────────
%\setbeamertemplate{footline}{%
%  \leavevmode%
%  \hbox{%
%    \begin{beamercolorbox}[wd=.35\paperwidth,ht=2.25ex,dp=1ex,center]{palette primary}%
%      \usebeamerfont{author in head/foot}G.~Gambosi
%    \end{beamercolorbox}%
%    \begin{beamercolorbox}[wd=.45\paperwidth,ht=2.25ex,dp=1ex,center]{palette secondary}%
%      \usebeamerfont{title in head/foot}Linear Regression
%    \end{beamercolorbox}%
%    \begin{beamercolorbox}[wd=.20\paperwidth,ht=2.25ex,dp=1ex,center]{palette tertiary}%
%      \usebeamerfont{date in head/foot}2025-2026 \quad \insertframenumber/\inserttotalframenumber
%    \end{beamercolorbox}%
%  }%
%  \vskip0pt%
%}
%\setbeamertemplate{navigation symbols}{}
%
%% ── Item style ───────────────────────────────────────────────────────────────
%\setbeamertemplate{itemize item}{\color{AccentBlue}\textbullet}
%\setbeamertemplate{itemize subitem}{\color{AccentBlue}$\circ$}
%\setbeamertemplate{enumerate item}{\color{MainBlue}\textbf{\insertenumlabel.}}

% ─────────────────────────────────────────────────────────────────────────────
%  TITLE
% ─────────────────────────────────────────────────────────────────────────────
\title[Linear Regression]{Linear Regression}
%\subtitle{Course of Machine Learning}
%\author{Giorgio Gambosi}
%\institute[Tor Vergata]{Master's Degree in Computer Science\\University of Rome ``Tor Vergata''}
%\date{A.Y. 2025–2026}

% ─────────────────────────────────────────────────────────────────────────────
%\begin{document}
% ─────────────────────────────────────────────────────────────────────────────

\begin{frame}
  \titlepage
\end{frame}

% ── Table of Contents ─────────────────────────────────────────────────────────
\begin{frame}{Outline}
  \tableofcontents
\end{frame}

% =============================================================================
\section{Linear Models and Basis Functions}
% =============================================================================

\begin{frame}{The Linear Predictor}
  \begin{block}{Definition}
    A \alert{linear model} predicts output as a weighted sum of input features plus a bias.
  \end{block}
  \vspace{.3em}
  For a $d$-dimensional input $\x = (x_1,\ldots,x_d)^T$:
  \[
    h(\x;\,w_1,\ldots,w_d,b) = w_1 x_1 + w_2 x_2 + \cdots + w_d x_d + b
  \]
  \begin{itemize}
    \item $w_i$ — weight for feature $x_i$.
    \item $b$ — \alert{bias}: shifts the decision surface from the origin.
  \end{itemize}
\end{frame}

\begin{frame}{Compact Vector Notation}
  Writing $\w = (w_1,\ldots,w_d)^T \in \Real^d$:
  \[
    h(\x;\w,b) = \w^T\x + b =
    \begin{pmatrix}w_1&w_2&\cdots&w_d\end{pmatrix}
    \begin{pmatrix}x_1\\x_2\\\vdots\\x_d\end{pmatrix} + b
  \]
  \vspace{.3em}
  \begin{alertblock}{Two senses of linearity}
    \begin{itemize}
      \item \textbf{Linearity in features} $\x$: hyperplane in input space.
      \item \textbf{Linearity in parameters} $\w, b$: convex objective $\Rightarrow$ global optimum is easy to find.
    \end{itemize}
  \end{alertblock}
\end{frame}

% ── Subsection: Feature Engineering ──────────────────────────────────────────
\subsection{Feature Engineering and Basis Functions}

\begin{frame}{Basis Function Expansion}
  \begin{block}{Idea}
    Map $\x\in\Real^d$ into a higher-dimensional space $\Real^m$ using \alert{basis functions}
    $\{\phi_1,\ldots,\phi_m\}$.
  \end{block}
  \[
    \Vphi(\x) = (\phi_1(\x),\ldots,\phi_m(\x))^T
  \]
  The model becomes:
  \[
    h(\x;\w,b) = \sum_{j=1}^m w_j\,\phi_j(\x) + b
  \]
  \begin{exampleblock}{Key observation}
    Basis functions \emph{do not} break linearity in parameters $\w$ —
    we still perform a linear combination, but of transformed inputs.
  \end{exampleblock}
\end{frame}

\begin{frame}{Types of Basis Functions}
  \begin{columns}[T]
  \column{.50\textwidth}
  \begin{block}{Polynomial (Global)}
    $\phi_j(x) = x^j$\\[2pt]
    {\small Affects prediction over the \emph{entire} input range.
    Can be unstable near boundaries.}
  \end{block}
  \vspace{.4em}
  \begin{block}{Gaussian (Local)}
    $\phi_j(x)=\exp\!\left(-\dfrac{(x-\mu_j)^2}{2s^2}\right)$\\[2pt]
    {\small Bell-shaped; significant only near centre $\mu_j$.}
  \end{block}
  \column{.50\textwidth}
  \begin{block}{Sigmoid (Local)}
    $\phi_j(x)=\sigma\!\left(\dfrac{x-\mu_j}{s}\right)=\dfrac{1}{1+e^{-(x-\mu_j)/s}}$\\[2pt]
    {\small Step-like transition centred at $\mu_j$.}
  \end{block}
  \vspace{.4em}
  \begin{block}{Hyperbolic Tangent (Local)}
    $\phi_j(x)=\tanh(x) = 2\sigma(x)-1$\\[2pt]
    {\small Maps to $[-1,1]$; common in neural networks.}
  \end{block}
  \end{columns}
\end{frame}

\begin{frame}{Basis Functions --- Illustration}
  \begin{center}
    \begin{tikzpicture}
      \begin{axis}[
        width=.82\textwidth, height=5.2cm,
        xlabel={$x$}, ylabel={$\phi(x)$},
        xmin=-3, xmax=3, ymin=-0.1, ymax=1.15,
        legend pos=north east,
        legend style={font=\small},
        grid=both, grid style={line width=.3pt,draw=gray!30},
        thick
      ]
      % Gaussian mu=0
      \addplot[MainBlue,domain=-3:3,samples=120]
        {exp(-x^2/0.5)};
      \addlegendentry{Gaussian $\mu{=}0$}
      % Gaussian mu=1.5
      \addplot[AccentBlue,domain=-3:3,samples=120,dashed]
        {exp(-(x-1.5)^2/0.5)};
      \addlegendentry{Gaussian $\mu{=}1.5$}
      % Sigmoid
      \addplot[AlertRed,domain=-3:3,samples=120]
        {1/(1+exp(-3*x))};
      \addlegendentry{Sigmoid $s{=}1/3$}
      % tanh
      \addplot[ExGreen,domain=-3:3,samples=120,densely dotted]
        {0.5*(tanh(2*x)+1)};
      \addlegendentry{tanh (shifted)}
      \end{axis}
    \end{tikzpicture}
  \end{center}
  \vspace{-.3em}
  {\small Different basis functions provide different types of local vs.\ global structure.}
\end{frame}

% ── Subsection: Design Matrix ─────────────────────────────────────────────────
\subsection{Design Matrix}

\begin{frame}{From Data to Design Matrix}
  \begin{columns}[T]
  \column{.48\textwidth}
  \begin{block}{Training data matrix $\X$}
    \[
      \X = \begin{pmatrix}
        x_{11}&x_{12}&\cdots&x_{1d}\\
        x_{21}&x_{22}&\cdots&x_{2d}\\
        \vdots&\vdots&\ddots&\vdots\\
        x_{n1}&x_{n2}&\cdots&x_{nd}
      \end{pmatrix}
    \]
    $n$ observations, $d$ features.
  \end{block}
  \column{.48\textwidth}
  \begin{block}{Design matrix $\VPhi$}
    \[
      \VPhi = \begin{pmatrix}
        \phi_1(\x_1)&\cdots&\phi_m(\x_1)\\
        \phi_1(\x_2)&\cdots&\phi_m(\x_2)\\
        \vdots&\ddots&\vdots\\
        \phi_1(\x_n)&\cdots&\phi_m(\x_n)
      \end{pmatrix}
    \]
    Each row: transformed sample.\\
    Each column: one basis function.
  \end{block}
  \end{columns}
  \vspace{.4em}
  \begin{itemize}
    \item All derivations for $\X$ carry over identically to $\VPhi$.
    \item Moving forward we write $(\X,\t)$ for clarity; $\VPhi$ is implied when using basis functions.
  \end{itemize}
\end{frame}

\begin{frame}{Running Example: Sinusoidal Regression}
  \begin{columns}[T]
  \column{.52\textwidth}
  \begin{exampleblock}{Setup}
    \begin{itemize}
      \item $n$ pairs $\{(x_i,t_i)\}$, $x\in[0,1]$.
      \item True signal: $t = \sin(2\pi x)$.
      \item Observations corrupted by Gaussian noise:
      \[ t_i = \sin(2\pi x_i) + \varepsilon_i, \quad \varepsilon_i\sim\mathcal{N}(0,\sigma^2) \]
    \end{itemize}
  \end{exampleblock}
  \column{.44\textwidth}
  \begin{center}
    \begin{tikzpicture}
      \begin{axis}[
        width=\textwidth, height=4.5cm,
        xlabel={$x$}, ylabel={$t$},
        xmin=0,xmax=1,ymin=-1.6,ymax=1.6,
        grid=both,grid style={line width=.2pt,draw=gray!25},
        tick label style={font=\tiny},
        label style={font=\small}
      ]
      % sine
      \addplot[MainBlue,thick,domain=0:1,samples=150]{sin(deg(2*pi*x))};
      % noisy points
      \addplot[only marks,mark=*,mark size=1.5pt,AccentBlue]
        coordinates{(0.06,0.30)(0.13,0.82)(0.20,1.05)(0.27,0.84)
                    (0.34,0.40)(0.41,-0.20)(0.48,-0.82)(0.55,-1.10)
                    (0.62,-0.88)(0.69,-0.41)(0.76,0.15)(0.83,0.72)
                    (0.90,0.98)(0.97,0.55)};
      \end{axis}
    \end{tikzpicture}
  \end{center}
  \end{columns}
  \vspace{.2em}
  Goal: \alert{generalize} — predict $t$ for unseen $\tilde{x}$ using only the noisy training points.
\end{frame}

\begin{frame}{Polynomial Regression as a Linear Model}
  \begin{block}{Polynomial hypothesis of degree $M$}
    \[
      h(x;\w,b) = \sum_{j=1}^M w_j\,x^j + b
    \]
    Basis functions: $\phi_j(x)=x^j$ — special case of the general framework.
  \end{block}
  \vspace{.3em}
  \begin{itemize}
    \item \textbf{Degree $M$ is the key model-selection hyperparameter.}
    \item Too low $M$: \alert{underfitting} — too rigid, misses oscillations.
    \item Too high $M$ (up to $M{+}1=n$): \alert{overfitting} — interpolates noise.
  \end{itemize}
\end{frame}

% =============================================================================
\section{Optimization: Analytical and Numerical}
% =============================================================================

\subsection{Closed-Form Solution (Normal Equations)}

\begin{frame}{Quadratic Loss and Convexity}
  \begin{block}{Empirical risk (sum of squared errors)}
    \[
      E(\ow) = \sum_{i=1}^n r_i(\ow)^2, \qquad r_i(\ow) = h(\x_i;\ow) - t_i
    \]
  \end{block}
  \begin{itemize}
    \item $E(\ow)$ is \textbf{convex}: sum of convex functions.
    \item Convexity $\Rightarrow$ any local minimum is the \alert{global minimum}.
    \item Derivative of $E$ is linear $\Rightarrow$ setting it to zero yields a \textbf{unique} solution.
  \end{itemize}
\end{frame}

\begin{frame}{Gradient and Normal Equations}
  Partial derivatives of $E(\ow)$:
  \[
    \frac{\partial E}{\partial w_k} = 2\sum_{i=1}^n r_i(\ow)\,x_{ik}, \qquad
    \frac{\partial E}{\partial b} = 2\sum_{i=1}^n r_i(\ow)
  \]
  Setting to zero $\Rightarrow$ linear system in matrix form: $\oX^T\oX\,\ow = \oX^T\t$\\[4pt]
  \textbf{Augmented matrix} $\oX$ (bias column prepended):
  \[
    \oX = \begin{pmatrix}1&x_{11}&\cdots&x_{1d}\\\vdots&\vdots&\ddots&\vdots\\1&x_{n1}&\cdots&x_{nd}\end{pmatrix}
  \]
  \begin{alertblock}{Normal equations — closed-form solution}
    \[
      \ow^* = (\oX^T\oX)^{-1}\oX^T\t
    \]
    Valid when $\oX^T\oX$ is non-singular (features not perfectly collinear).
  \end{alertblock}
\end{frame}

\subsection{Gradient Descent}

\begin{frame}{Gradient Descent — Algorithm}
  \begin{block}{When to use}
    Closed-form inversion is $O(d^3)$: too costly for large $d$ or $n$.
    Gradient descent iterates cheaply.
  \end{block}
  \vspace{.3em}
  \begin{enumerate}
    \item \textbf{Initialization:} choose $\ow^{(0)}$ (e.g.\ zero or random).
    \item \textbf{Descent step:} move in direction of \alert{steepest descent}:
    \[
      \ow^{(s)} := \ow^{(s-1)} - \eta\,\nabla E(\ow)\Big|_{\ow^{(s-1)}}
    \]
    \item \textbf{Component-wise update:}
    \[
      w_k^{(s)} := w_k^{(s-1)} - 2\eta\sum_{i=1}^n r_i(\ow^{(s-1)})\,x_{ik}
    \]
  \end{enumerate}
  \begin{itemize}
    \item $\eta$ = \alert{learning rate} (step size).
    \item Update proportional to residual $r_i$: large error $\Rightarrow$ large correction.
  \end{itemize}
\end{frame}

\begin{frame}{Matrix Form of the Update}
  \begin{block}{Compact update rule}
    \[
      \ow^{(s)} := \ow^{(s-1)} - 2\eta\sum_{i=1}^n r_i(\ow^{(s-1)})\,\ox_i
    \]
    where $\ox_i$ is the $i$-th augmented row of $\oX$.
  \end{block}
  \vspace{.4em}
  \begin{columns}[T]
    \column{.48\textwidth}
    \begin{alertblock}{$\eta$ too small}
      Convergence is very slow; many iterations required.
    \end{alertblock}
    \column{.48\textwidth}
    \begin{alertblock}{$\eta$ too large}
      Overshoots the minimum; may diverge.
    \end{alertblock}
  \end{columns}
\end{frame}

% =============================================================================
\section{Model Selection and Overfitting}
% =============================================================================

\subsection{Polynomial Degree and Generalization}

\begin{frame}{Overfitting vs.\ Underfitting}
  \begin{columns}[T]
  \column{.50\textwidth}
  \begin{block}{The trade-off}
    \begin{itemize}
      \item As $M\uparrow$: training error $\downarrow$ monotonically.
      \item Test error is \textbf{non-monotone}: decreases then rises.
      \item The model eventually memorises \emph{noise}, not signal.
    \end{itemize}
  \end{block}
  \column{.46\textwidth}
  \begin{center}
    \begin{tikzpicture}
      \begin{axis}[
        width=\textwidth,height=4.4cm,
        xlabel={Polynomial degree $M$},
        ylabel={$E_{\mathrm{RMS}}$},
        xmin=0,xmax=9,ymin=0,ymax=1,
        xtick={0,1,2,3,4,5,6,7,8,9},
        legend pos=north east,
        legend style={font=\tiny},
        grid=both,grid style={line width=.2pt,draw=gray!25},
        thick
      ]
      \addplot[MainBlue,mark=square*,mark size=1.5pt]
        coordinates{(0,.8)(1,.7)(2,.5)(3,.28)(4,.22)(5,.18)(6,.12)(7,.08)(8,.04)(9,.01)};
      \addlegendentry{Train}
      \addplot[AlertRed,mark=*,mark size=1.5pt]
        coordinates{(0,.83)(1,.73)(2,.53)(3,.30)(4,.26)(5,.33)(6,.52)(7,.70)(8,.85)(9,.95)};
      \addlegendentry{Test}
      \end{axis}
    \end{tikzpicture}
  \end{center}
  \end{columns}
  \vspace{.2em}
  \begin{itemize}
    \item RMS error: $E_{\mathrm{RMS}}(\w^*,\mathcal{T}) = \sqrt{\frac{1}{|\mathcal{T}|}\sum_i(h(x_i;\w^*)-t_i)^2}$
    \item Low test error = successful \alert{generalisation}.
  \end{itemize}
\end{frame}

\begin{frame}{Polynomial Fits at Different Degrees}
  \begin{center}
  \begin{tikzpicture}
    \begin{axis}[
      name=ax1, width=.4\textwidth, height=4cm,
      title={$M=0$: Underfitting},
      title style={font=\small},
      xmin=0,xmax=1,ymin=-1.5,ymax=1.5,
      grid=both,grid style={line width=.2pt,draw=gray!25},
      tick label style={font=\tiny},
      xlabel style={font=\tiny}, ylabel style={font=\tiny}
    ]
    \addplot[MainBlue,thick,domain=0:1]{sin(deg(2*pi*x))};
    \addplot[AlertRed,thick,domain=0:1]{0.05};
    \addplot[only marks,mark=*,mark size=1.5pt,AccentBlue]
      coordinates{(0.1,0.55)(0.2,1.0)(0.35,0.6)(0.5,-0.9)(0.65,-0.7)(0.8,0.55)(0.9,0.9)};
    \end{axis}
    \begin{axis}[
      name=ax3, at={(ax1.east)}, xshift=1.2cm,
      width=.4\textwidth, height=4cm,
      title={$M=3$: Good fit},
      title style={font=\small},
      xmin=0,xmax=1,ymin=-1.5,ymax=1.5,
      grid=both,grid style={line width=.2pt,draw=gray!25},
      tick label style={font=\tiny},
      xlabel style={font=\tiny}, ylabel style={font=\tiny}
    ]
    \addplot[MainBlue,thick,domain=0:1]{sin(deg(2*pi*x))};
    \addplot[AlertRed,thick,domain=0:1,samples=100]{sin(deg(2*pi*x))*0.92-0.05*x};
    \addplot[only marks,mark=*,mark size=1.5pt,AccentBlue]
      coordinates{(0.1,0.55)(0.2,1.0)(0.35,0.6)(0.5,-0.9)(0.65,-0.7)(0.8,0.55)(0.9,0.9)};
    \end{axis}
  \end{tikzpicture}
  \end{center}
  \vspace{.1em}
  \begin{itemize}
    \item $M{=}0$: constant prediction (flat line) — very high bias.
    \item $M{=}3$: captures sinusoidal trend well.
    \item $M{=}9$ (with $n{=}10$ points): passes through every point — zero training error but poor test error.
  \end{itemize}
\end{frame}

\subsection{Training Set Size}

\begin{frame}{Effect of Training Set Size}
  \begin{block}{Key observation}
    For a fixed model complexity $M$, overfitting \emph{decreases} as $n$ increases.
  \end{block}
  \begin{columns}[T]
  \column{.48\textwidth}
  \begin{alertblock}{Small dataset}
    \begin{itemize}
      \item Few constraints on the polynomial.
      \item High-degree curve wiggles freely.
      \item \alert{High overfitting risk.}
    \end{itemize}
  \end{alertblock}
  \column{.48\textwidth}
  \begin{exampleblock}{Large dataset}
    \begin{itemize}
      \item Many constraints force smooth fit.
      \item Harder to fit individual noise samples.
      \item \alert{Reduced overfitting.}
    \end{itemize}
  \end{exampleblock}
  \end{columns}
  \vspace{.4em}
  \begin{itemize}
    \item Rule of thumb: \textbf{larger $n$ tolerates larger $M$}.
    \item ``Big Data'' unlocks more complex models.
    \item When more data is unavailable $\Rightarrow$ use \alert{regularisation}.
  \end{itemize}
\end{frame}

% =============================================================================
\section{Regularisation}
% =============================================================================

\subsection{Regularised Cost Function}

\begin{frame}{Adding a Regularisation Term}
  \begin{block}{Total cost function}
    \[
      E_{\mathrm{total}}(\ow) = E_D(\ow) + \lambda\,E_W(\ow)
    \]
  \end{block}
  \begin{itemize}
    \item $E_D(\ow)$ — \alert{data term} (empirical risk): measures fit to observations.
    \item $E_W(\ow)$ — \alert{regularisation term}: penalises extreme parameter values.
    \item $\lambda$ — \alert{regularisation coefficient}:
      \begin{itemize}
        \item Small $\lambda$: prioritise fit $\Rightarrow$ risk overfitting.
        \item Large $\lambda$: prioritise smoothness $\Rightarrow$ risk underfitting.
      \end{itemize}
  \end{itemize}
\end{frame}

\subsection{Ridge Regression ($q=2$)}

\begin{frame}{Ridge Regression (Tikhonov Regularisation)}
  \begin{block}{Quadratic regularisation}
    \[
      E_W(\ow) = \frac{1}{2}\ow^T\ow = \frac{1}{2}\sum_{i=1}^m w_i^2 + \frac{b^2}{2}
    \]
  \end{block}
  Full objective:
  \[
    E(\ow) = \frac{1}{2}\boldsymbol{r}(\ow)^T\boldsymbol{r}(\ow) + \frac{\lambda}{2}\ow^T\ow
  \]
  \begin{alertblock}{Regularised normal equations}
    \[
      \ow = \bigl(\lambda\I + \oX^T\oX\bigr)^{-1}\oX^T\t
    \]
    $\lambda\I$ strengthens the diagonal $\Rightarrow$ always well-conditioned.
  \end{alertblock}
  \begin{itemize}
    \item Closed-form solution retained.
    \item Shrinks all weights toward zero, but rarely to \emph{exactly} zero.
  \end{itemize}
\end{frame}

\subsection{Lasso ($q=1$) and Sparse Models}

\begin{frame}{Generalised Regularisation and Lasso}
  General form with exponent $q$:
  \[
    E(\ow) = \frac{1}{2}\sum_i r_i(\ow)^2 + \frac{\lambda}{2}\sum_j |w_j|^q
  \]
  \begin{columns}[T]
  \column{.48\textwidth}
  \begin{block}{Ridge ($q=2$)}
    \begin{itemize}
      \item Constraint: $\sum w_j^2\leq\eta$ — \alert{hypersphere}.
      \item Smooth boundary.
      \item All weights shrink toward 0; none exactly 0.
      \item Retains all features.
    \end{itemize}
  \end{block}
  \column{.48\textwidth}
  \begin{block}{Lasso ($q=1$)}
    \begin{itemize}
      \item Constraint: $\sum|w_j|\leq\eta$ — \alert{hyperdiamond}.
      \item Sharp corners on axes.
      \item Error contours hit corners $\Rightarrow$ weights set to \textbf{exactly 0}.
      \item Built-in \alert{feature selection}.
    \end{itemize}
  \end{block}
  \end{columns}
\end{frame}

\begin{frame}{Geometric Interpretation: Ridge vs.\ Lasso}
  \begin{center}
  \begin{tikzpicture}[scale=1.0]
    % ---- Ridge circle (left) ----
    \draw[MainBlue,thick,fill=torgray] (-3.2,0) circle(1.2);
    \draw[AlertRed,thick] (-1.7,0.6) ellipse(0.9 and 0.5);
    \draw[AlertRed,dashed] (-1.7,0.6) ellipse(1.6 and 0.9);
    \draw[AlertRed,densely dotted] (-1.7,0.6) ellipse(2.4 and 1.35);
    \draw[->,gray!60] (-5,0)--(-1.4,0) node[right,font=\tiny]{$w_1$};
    \draw[->,gray!60] (-3.2,-1.6)--(-3.2,1.6) node[above,font=\tiny]{$w_2$};
    \fill[AlertRed] (-2.36,0.87) circle(2.5pt);
    \node[font=\scriptsize,right] at (-2.36,0.87) {$\ow^*$};
    \node[font=\small,below] at (-3.2,-1.85) {Ridge ($q=2$)};
    % ---- Lasso diamond (right) ----
    \draw[MainBlue,thick,fill=torgray] (3.2,1.2)--(4.4,0)--(3.2,-1.2)--(2.0,0)--cycle;
    \draw[AlertRed,thick] (4.7,0.6) ellipse(0.9 and 0.5);
    \draw[AlertRed,dashed] (4.7,0.6) ellipse(1.6 and 0.9);
    \draw[AlertRed,densely dotted] (4.7,0.6) ellipse(2.4 and 1.35);
    \draw[->,gray!60] (1.4,0)--(5.0,0) node[right,font=\tiny]{$w_1$};
    \draw[->,gray!60] (3.2,-1.6)--(3.2,1.6) node[above,font=\tiny]{$w_2$};
    \fill[AlertRed] (4.4,0) circle(2.5pt);
    \node[font=\scriptsize,above right] at (4.4,0) {$\ow^*$\;($w_2{=}0$)};
    \node[font=\small,below] at (3.2,-1.85) {Lasso ($q=1$)};
  \end{tikzpicture}
  \end{center}
  \vspace{-.3em}
  {\small Lasso corners lie on axes $\Rightarrow$ error contours likely hit a corner $\Rightarrow$ sparse solution.}
\end{frame}

\begin{frame}{Effect of $\lambda$ on Model Fit}
  \begin{columns}[T]
  \column{.33\textwidth}
  \begin{alertblock}{Small $\lambda$}
    \begin{itemize}
      \item Penalty negligible.
      \item Low training error.
      \item High test error (\alert{overfitting}).
    \end{itemize}
  \end{alertblock}
  \column{.33\textwidth}
  \begin{exampleblock}{Intermediate $\lambda$}
    \begin{itemize}
      \item Penalty balances fit.
      \item Moderate training error.
      \item \alert{Minimum test error.}
    \end{itemize}
  \end{exampleblock}
  \column{.33\textwidth}
  \begin{alertblock}{Large $\lambda$}
    \begin{itemize}
      \item Penalty dominates.
      \item Weights heavily suppressed.
      \item High train \& test error (\alert{underfitting}).
    \end{itemize}
  \end{alertblock}
  \end{columns}
  \vspace{.5em}
  \begin{center}
  \begin{tikzpicture}
    \begin{axis}[
      width=.7\textwidth, height=3.8cm,
      xlabel={$\ln\lambda$},
      ylabel={$E_{\mathrm{RMS}}$},
      xmin=-5,xmax=3, ymin=0, ymax=1,
      legend pos=north west, legend style={font=\small},
      grid=both,grid style={line width=.2pt,draw=gray!25},thick
    ]
    \addplot[MainBlue,mark=none,domain=-5:3,samples=80]
      {0.05 + 0.01*exp(0.9*(x+2))};
    \addlegendentry{Train}
    \addplot[AlertRed,mark=none,domain=-5:3,samples=80]
      {0.1 + 0.08*exp(-1.5*(x+1)) + 0.07*exp(0.8*(x+1))};
    \addlegendentry{Test}
    \end{axis}
  \end{tikzpicture}\\
  {\small Error on training (blue) and test (red) sets as a function of $\ln\lambda$.}
  \end{center}
\end{frame}

% ── Bias-Variance ──────────────────────────────────────────────────────────────
\subsection{Bias-Variance Trade-off}

\begin{frame}{Bias-Variance Decomposition}
  \begin{block}{Setup}
    $L=100$ independent datasets; $n=25$ points each; $m=24$ Gaussian basis functions.
    Three regularisation regimes:
  \end{block}
  \vspace{.3em}
  \begin{columns}[T]
  \column{.32\textwidth}
  \begin{alertblock}{High $\lambda$ ($\ln\lambda=2.6$)}
    \begin{itemize}
      \item Models very similar.
      \item \alert{Low variance.}
      \item Average far from sine.
      \item \alert{High bias.}
    \end{itemize}
  \end{alertblock}
  \column{.32\textwidth}
  \begin{exampleblock}{Medium $\lambda$ ($\ln\lambda=-0.31$)}
    \begin{itemize}
      \item Good balance.
      \item Follows sinusoidal trend.
      \item Average close to $\sin(2\pi x)$.
    \end{itemize}
  \end{exampleblock}
  \column{.32\textwidth}
  \begin{alertblock}{Low $\lambda$ ($\ln\lambda=-2.4$)}
    \begin{itemize}
      \item Models wildly different.
      \item \alert{High variance.}
      \item Average close to sine.
      \item \alert{Low bias.}
    \end{itemize}
  \end{alertblock}
  \end{columns}
\end{frame}

\begin{frame}{Bias-Variance Trade-off — Graph}
  \begin{columns}[T]
  \column{.50\textwidth}
  \begin{center}
  \begin{tikzpicture}
    \begin{axis}[
      width=\textwidth, height=5.5cm,
      xlabel={$\ln\lambda$},
      ylabel={Error},
      xmin=-3,xmax=3,ymin=0,ymax=1.2,
      legend pos=north east, legend style={font=\small},
      grid=both,grid style={line width=.2pt,draw=gray!25},thick
    ]
    % Bias^2: increases with lambda
    \addplot[MainBlue,domain=-3:3,samples=80]
      {0.02 + 0.12*exp(0.7*(x+1))};
    \addlegendentry{$(\text{Bias})^2$}
    % Variance: decreases with lambda
    \addplot[AlertRed,domain=-3:3,samples=80]
      {0.6*exp(-0.8*(x+2))};
    \addlegendentry{Variance}
    % Total
    \addplot[ExGreen,thick,domain=-3:3,samples=80]
      {0.02+0.12*exp(0.7*(x+1)) + 0.6*exp(-0.8*(x+2))};
    \addlegendentry{Total}
    \end{axis}
  \end{tikzpicture}
  \end{center}
  \column{.46\textwidth}
  \begin{block}{Key insight}
    \begin{itemize}
      \item $\lambda\uparrow$: model simpler $\Rightarrow$ \alert{bias $\uparrow$, variance $\downarrow$}.
      \item $\lambda\downarrow$: model complex $\Rightarrow$ \alert{bias $\downarrow$, variance $\uparrow$}.
      \item Optimal $\lambda$: minimises \textbf{total} error $= (\text{bias})^2 + \text{variance}$.
    \end{itemize}
  \end{block}
  \vspace{.3em}
  {\small Overfitting = high variance. Underfitting = high bias.}
  \end{columns}
\end{frame}

% =============================================================================
\section{Probabilistic Model for Regression}
% =============================================================================

\subsection{Gaussian Noise Model}

\begin{frame}{Probabilistic Framework}
  \begin{block}{Motivation}
    Move beyond point fitting; model \alert{uncertainty} in data via probability.
  \end{block}
  \vspace{.3em}
  \textbf{Assumption:} target $t$ is Gaussian around hypothesis $h(\x;\ow)$:
  \[
    p(t|\x;\ow,\beta) = \mathcal{N}\bigl(t\mid h(\x;\ow),\,\beta^{-1}\bigr)
  \]
  \begin{itemize}
    \item $h(\x;\ow)$ = \alert{mean}: our prediction.
    \item $\beta = 1/\sigma^2$ = \alert{precision}: inverse noise variance.
    \item High $\beta$ $\Rightarrow$ tight cluster around prediction.
  \end{itemize}
  \vspace{.3em}
  \begin{center}
  \begin{tikzpicture}
    \begin{axis}[
      width=.7\textwidth, height=3.8cm,
      xlabel={$x$}, ylabel={$t$},
      xmin=0,xmax=1,ymin=-1.8,ymax=1.8,
      grid=both,grid style={line width=.2pt,draw=gray!25},
      tick label style={font=\tiny},thick
    ]
    \addplot[MainBlue,domain=0:1,samples=100]{sin(deg(2*pi*x))};
    \addplot[AccentBlue,name path=upA,domain=0:1,samples=100]{sin(deg(2*pi*x))+0.4};
    \addplot[AccentBlue,name path=loA,domain=0:1,samples=100]{sin(deg(2*pi*x))-0.4};
    \addplot[AccentBlue,fill opacity=0.15,draw=none] fill between[of=upA and loA];
    \addplot[only marks,mark=*,mark size=1.5pt,AlertRed]
      coordinates{(0.1,0.5)(0.2,1.1)(0.35,0.7)(0.5,-0.8)(0.65,-0.65)(0.8,0.5)};
    \end{axis}
  \end{tikzpicture}
  \end{center}
\end{frame}

% ── MLE ────────────────────────────────────────────────────────────────────────
\subsection{Maximum Likelihood Estimation (MLE)}

\begin{frame}{Likelihood Function}
  i.i.d.\ observations $\Rightarrow$ likelihood is the product:
  \[
    L(\ow,\beta|\X,\t) = \prod_{i=1}^n \mathcal{N}\bigl(t_i\mid h(\x_i;\ow),\beta^{-1}\bigr)
    = \prod_{i=1}^n \frac{\sqrt{\beta}}{\sqrt{2\pi}}\,e^{-\frac{\beta}{2}r_i(\ow)^2}
  \]
  \textbf{Log-likelihood} (turns product into sum):
  \[
    l(\ow,\beta) = -\frac{\beta}{2}\sum_{i=1}^n r_i(\ow)^2 + \frac{n}{2}\log\beta + \text{const}
  \]
\end{frame}

\begin{frame}{MLE Estimates}
  \begin{columns}[T]
  \column{.50\textwidth}
  \begin{block}{Optimal weights $\ow_{ML}$}
    Maximising $l$ over $\ow$ $\equiv$ minimising $\sum r_i^2$.\\[4pt]
    \alert{MLE = Least Squares} under Gaussian noise.\\
    $\ow_{ML} = (\oX^T\oX)^{-1}\oX^T\t$
  \end{block}
  \column{.46\textwidth}
  \begin{block}{Optimal precision $\beta_{ML}$}
    Setting $\partial l/\partial\beta = 0$:
    \[
      \beta_{ML}^{-1} = \frac{1}{n}\sum_{i=1}^n r_i(\ow_{ML})^2
    \]
    Average squared residual.
  \end{block}
  \end{columns}
  \vspace{.4em}
  \begin{exampleblock}{Predictive distribution (MLE)}
    For a new input $\x$:
    \[
      p(t|\x;\ow_{ML},\beta_{ML}) = \mathcal{N}\bigl(t\mid h(\x;\ow_{ML}),\,\beta_{ML}^{-1}\bigr)
    \]
  \end{exampleblock}
\end{frame}

% =============================================================================
\section{Bayesian Linear Regression}
% =============================================================================

\subsection{Prior and Posterior}

\begin{frame}{From Frequentist to Bayesian}
  \begin{alertblock}{MLE limitation}
    Treats $\ow$ as a fixed unknown $\Rightarrow$ prone to overfitting on small data.
  \end{alertblock}
  \vspace{.3em}
  \begin{block}{Bayesian view}
    Treat $\ow$ as a \alert{random variable} with a prior distribution.
  \end{block}
  \vspace{.3em}
  \textbf{Zero-mean Gaussian prior} (hyperparameter $\alpha$):
  \[
    p(\ow|\alpha) = \mathcal{N}(\ow;\,\Zero,\,\alpha^{-1}\I) = \left(\frac{\alpha}{2\pi}\right)^{m/2}\!
    e^{-\frac{\alpha}{2}\ow^T\ow}
  \]
  \begin{itemize}
    \item Encodes belief that weights are small (close to zero).
    \item $\alpha$ controls the \alert{strength of the prior}.
  \end{itemize}
\end{frame}

\begin{frame}{Conjugate Prior and Posterior}
  By Bayes' rule:
  \[
    p(\ow|\X,\t;\alpha,\beta) \propto p(\t|\X,\ow;\beta)\,p(\ow|\alpha)
  \]
  \begin{block}{Gaussian prior is conjugate $\Rightarrow$ posterior is Gaussian}
    \begin{align*}
      \VSigma_p &= \bigl(\alpha\I + \beta\oX^T\oX\bigr)^{-1}\\[4pt]
      \mathbf{m}_p &= \beta\,\VSigma_p\,\oX^T\t
    \end{align*}
  \end{block}
  \begin{exampleblock}{Connection to Ridge Regression}
    $\mathbf{m}_p$ is identical to the Ridge solution with $\lambda = \alpha/\beta$.\\
    \alert{Regularisation = Gaussian prior on weights.}
  \end{exampleblock}
\end{frame}

\subsection{Sequential Bayesian Learning}

\begin{frame}{Sequential Learning}
  \begin{block}{Core idea}
    The posterior after observing $\mathcal{T}_{1},\ldots,\mathcal{T}_{k-1}$ becomes the \alert{prior} for step $k$.
  \end{block}
  \vspace{.3em}
  Update chain:
  \begin{align*}
    p(\ow|\mathcal{T}_1) &\propto p(\mathcal{T}_1|\ow)\,p(\ow)\\
    p(\ow|\mathcal{T}_1,\ldots,\mathcal{T}_k) &\propto p(\mathcal{T}_k|\ow)\,p(\ow|\mathcal{T}_1,\ldots,\mathcal{T}_{k-1})
  \end{align*}
  \begin{itemize}
    \item \textbf{Online learning}: processes data one sample (or batch) at a time.
    \item Each new observation narrows the posterior.
    \item Full-batch result is mathematically identical to sequential result.
  \end{itemize}
\end{frame}

\begin{frame}{Sequential Learning — Visual Example (1/2)}
  \begin{columns}[T]
  \column{.48\textwidth}
  \begin{block}{Setup}
    Model: $h(x;b,w_1)=w_1 x+b$\\
    True: $y=-0.3+0.5x$ + noise $\sigma=0.2$\\
    Prior: $\mathcal{N}(\Zero, 0.04\I)$
  \end{block}
  \vspace{.3em}
  \textbf{Initial state (no data):}\\
  Prior is a broad circle centred at $(0,0)$.\\
  Lines sampled from prior are random in orientation.
  \column{.48\textwidth}
  \begin{center}
  \begin{tikzpicture}
    \begin{axis}[
      width=\textwidth,height=5cm,
      xlabel={$b$},ylabel={$w_1$},
      xmin=-1,xmax=1,ymin=-1,ymax=1,
      title={Prior distribution},
      title style={font=\small},
      grid=both,grid style={line width=.2pt,draw=gray!25},
      tick label style={font=\tiny}
    ]
    % broad prior
    \draw[MainBlue,thick,fill=AccentBlue,fill opacity=0.2]
      (axis cs:0,0) circle[radius=30pt];
    \addplot[only marks,mark=*,mark size=2pt,MainBlue]
      coordinates{(-0.3,0.5)};
    \node[font=\tiny,right,MainBlue] at (axis cs:-0.3,0.55) {true};
    \end{axis}
  \end{tikzpicture}
  \end{center}
  \end{columns}
\end{frame}

\begin{frame}{Sequential Learning — Visual Example (2/2)}
  \begin{itemize}
    \item After \textbf{1 observation}: posterior elongated into a stripe — many $(b,w_1)$ pairs consistent with one point.
    \item After \textbf{2 observations}: intersection of two stripes; sampled lines pass near both points.
    \item After \textbf{n observations}: posterior concentrates tightly around $(a_0{=}-0.3,\,a_1{=}0.5)$.
  \end{itemize}
  \vspace{.3em}
  \begin{center}
  \begin{tikzpicture}
    \begin{axis}[
      width=.55\textwidth, height=4cm,
      xlabel={$b$}, ylabel={$w_1$},
      xmin=-1,xmax=1,ymin=-1,ymax=1,
      title={Posterior after many observations},
      title style={font=\small},
      grid=both,grid style={line width=.2pt,draw=gray!25},
      tick label style={font=\tiny}
    ]
    \draw[MainBlue,thick,fill=AccentBlue,fill opacity=0.3]
      (axis cs:-0.3,0.5) circle[radius=6pt];
    \addplot[only marks,mark=star,mark size=3pt,AlertRed]
      coordinates{(-0.3,0.5)};
    \node[font=\tiny,right,AlertRed] at (axis cs:-0.28,0.52) {$(a_0,a_1)$};
    \end{axis}
  \end{tikzpicture}
  \end{center}
  Posterior variance shrinks monotonically: $\VSigma_{n+1}\preceq\VSigma_n$.
\end{frame}

\begin{frame}{Why Does Posterior Variance Decrease?}
  Posterior precision update when a new point $(\x_{n+1},t_{n+1})$ arrives:
  \[
    \VSigma_{n+1}^{-1} = \VSigma_n^{-1} + \beta\,\ox_{n+1}\ox_{n+1}^T
  \]
  \begin{itemize}
    \item $\ox_{n+1}\ox_{n+1}^T$ is \textbf{positive semi-definite} $\Rightarrow$ precision increases.
    \item Larger precision $\Leftrightarrow$ smaller covariance.
    \item Via Sherman-Morrison:
    \[
      \VSigma_{n+1} = \VSigma_n - \frac{\VSigma_n\ox\ox^T\VSigma_n}{\beta^{-1}+\ox^T\VSigma_n\ox}
    \]
    \item $\Delta\VSigma = \VSigma_n - \VSigma_{n+1}$ is positive semi-definite: \alert{variance never increases}.
  \end{itemize}
\end{frame}

\subsection{Maximum A Posteriori (MAP)}

\begin{frame}{MAP Estimation}
  \begin{block}{Definition}
    $\ow_{MAP}$ = \alert{mode} of the posterior — the most probable parameter vector:
    \[
      \ow_{MAP} = \argmax{\ow}\,p(\ow|\X,\t;\alpha,\beta)
    \]
  \end{block}
  Equivalent minimisation (drop constant evidence term):
  \[
    \ow_{MAP} = \argmin{\ow}\left[-\log p(\t|\X,\ow;\beta) - \log p(\ow;\alpha)\right]
  \]
  Expanding with Gaussian assumptions:
  \[
    \frac{\beta}{2}\sum_i r_i(\ow)^2 + \frac{\alpha}{2}\sum_j\overline{w}_j^2
    \;\propto\;
    \frac{1}{2}\sum_i r_i(\w)^2 + \frac{\lambda}{2}\|\ow\|^2
  \]
  \begin{alertblock}{MAP = Ridge Regression with $\lambda = \alpha/\beta$}
    Regularisation is a \alert{probabilistic} consequence of a Gaussian prior.
  \end{alertblock}
\end{frame}

\begin{frame}{MAP Predictive Distribution}
  With the MAP estimate, predict via:
  \[
    p(t|\x,\ow_{MAP};\beta) = \mathcal{N}\bigl(t\mid h(\x;\ow_{MAP}),\,\beta^{-1}\bigr)
  \]
  \begin{itemize}
    \item Same form as MLE predictive distribution.
    \item Uses $\ow_{MAP}$ (regularised) instead of $\ow_{ML}$.
    \item Still a \textbf{point estimate} of $\ow$ — ignores parameter uncertainty.
    \item Fully Bayesian approach: integrate over all $\ow$.
  \end{itemize}
\end{frame}

% =============================================================================
\section{Approaches to Prediction}
% =============================================================================

\begin{frame}{Three Approaches to Prediction in Linear Regression}
  \begin{block}{Classical (Least Squares / MLE)}
    \begin{itemize}
      \item Point estimate $\ow_{LS}$ via minimising quadratic cost (or MLE under Gaussian noise).
      \item Prediction: $t = \ox^T\ow_{LS}$.
      \item Regularisation can be added to reduce overfitting.
    \end{itemize}
  \end{block}
\end{frame}

\begin{frame}{Three Approaches (continued)}
  \begin{block}{Bayesian Point Estimation (MAP)}
    \begin{itemize}
      \item Derive posterior $p(\ow|\X,\t;\alpha,\beta)$; extract mode $\ow_{MAP}$.
      \item Equivalent to Ridge: $\ow_{MAP}=\ow_{LS}$ when $\lambda=\alpha/\beta$.
      \item Prediction: Gaussian $\mathcal{N}(t\mid\ox^T\ow_{MAP},\beta^{-1})$.
    \end{itemize}
  \end{block}
  \vspace{.3em}
  \begin{block}{Fully Bayesian}
    \begin{itemize}
      \item Goal: \alert{predictive distribution} $p(t|\x,\X,\t)$ marginalising over $\ow$:
      \[
        p(t|\x,\X,\t;\alpha,\beta) = \int p(t|\x,\ow;\beta)\,p(\ow|\X,\t;\alpha,\beta)\,d\ow
      \]
      \item Convolution of two Gaussians $\Rightarrow$ \textbf{Gaussian} result.
      \item Accounts for parameter uncertainty — not just noise.
    \end{itemize}
  \end{block}
\end{frame}

% =============================================================================
\section{Fully Bayesian Predictive Distribution}
% =============================================================================

\subsection{Predictive Mean and Variance}

\begin{frame}{The Fully Bayesian Predictive Distribution}
  \begin{block}{Result}
    \[
      p(t|\x,\X,\t;\alpha,\beta) = \mathcal{N}(t;\;m(\x),\;\sigma^2(\x))
    \]
    with:
    \[
      m(\x) = \beta\,\ox^T\VSigma\,\oX^T\t, \qquad
      \sigma^2(\x) = \frac{1}{\beta} + \ox^T\VSigma\,\ox
    \]
    where $\VSigma = (\alpha\I+\beta\oX^T\oX)^{-1}$.
  \end{block}
  \vspace{.3em}
  \textbf{Two components of predictive variance:}
  \[
    \sigma^2(\x) = \underbrace{\frac{1}{\beta}}_{\text{aleatoric (noise)}} + \underbrace{\ox^T\VSigma\,\ox}_{\text{epistemic (model uncertainty)}}
  \]
\end{frame}

\begin{frame}{Aleatoric vs.\ Epistemic Uncertainty}
  \begin{columns}[T]
  \column{.50\textwidth}
  \begin{block}{Aleatoric uncertainty $1/\beta$}
    \begin{itemize}
      \item Intrinsic noise in the data.
      \item Irreducible: cannot be reduced even with more data.
      \item Represents measurement errors, unobserved variables.
    \end{itemize}
  \end{block}
  \column{.46\textwidth}
  \begin{block}{Epistemic uncertainty $\ox^T\VSigma\ox$}
    \begin{itemize}
      \item Our ignorance of the true $\ow$.
      \item \alert{Reducible}: decreases as training data grows.
      \item Small near training points; large in data-void regions.
    \end{itemize}
  \end{block}
  \end{columns}
  \vspace{.4em}
  \begin{exampleblock}{Geometric insight}
    At a test point close to training data: $\ox$ aligns with well-informed directions
    $\Rightarrow$ $\ox^T\VSigma\ox$ is small.\\
    In data-void gaps: $\ox$ points in under-constrained directions $\Rightarrow$ large uncertainty.
  \end{exampleblock}
\end{frame}

\subsection{Evolution of Predictive Distribution}

\begin{frame}{Bayesian Predictive Distribution — $n=1$}
  \begin{exampleblock}{Experimental setup}
    Function: $f(x)=\sin(2\pi x)$; model: $m=9$ Gaussian basis functions; observations with Gaussian noise.
  \end{exampleblock}
  \vspace{.2em}
  \begin{columns}[T]
  \column{.50\textwidth}
  \begin{center}
  \begin{tikzpicture}
    \begin{axis}[
      width=\textwidth,height=5.2cm,
      xlabel={$x$}, ylabel={$t$},
      xmin=0,xmax=1,ymin=-2.5,ymax=2.5,
      grid=both,grid style={line width=.2pt,draw=gray!25},
      tick label style={font=\tiny},thick,
      title={$n=1$: Predictive distribution},
      title style={font=\small}
    ]
    \addplot[MainBlue,domain=0:1,samples=100]{sin(deg(2*pi*x))};
    % wide uncertainty
    \addplot[AlertRed,name path=mu,domain=0:1,samples=100]{0.1*(x-0.5)};
    \addplot[AccentBlue,name path=up,domain=0:1,samples=100]{0.1*(x-0.5)+1.8};
    \addplot[AccentBlue,name path=lo,domain=0:1,samples=100]{0.1*(x-0.5)-1.8};
    \addplot[AccentBlue,fill opacity=0.2,draw=none] fill between[of=up and lo];
    \addplot[only marks,mark=*,mark size=2pt,black] coordinates{(0.35,0.6)};
    \end{axis}
  \end{tikzpicture}
  \end{center}
  \column{.46\textwidth}
  \textbf{With only 1 data point:}
  \begin{itemize}
    \item Predictive mean (red) is nearly flat.
    \item Pink band extremely wide: \alert{high uncertainty} everywhere.
    \item Many different functions are equally plausible.
  \end{itemize}
  \end{columns}
\end{frame}

\begin{frame}{Bayesian Predictive Distribution — $n=4$}
  \begin{columns}[T]
  \column{.50\textwidth}
  \begin{center}
  \begin{tikzpicture}
    \begin{axis}[
      width=\textwidth,height=5.2cm,
      xlabel={$x$}, ylabel={$t$},
      xmin=0,xmax=1,ymin=-2.2,ymax=2.2,
      grid=both,grid style={line width=.2pt,draw=gray!25},
      tick label style={font=\tiny},thick,
      title={$n=4$: Predictive distribution},
      title style={font=\small}
    ]
    \addplot[MainBlue,domain=0:1,samples=100]{sin(deg(2*pi*x))};
    \addplot[AlertRed,domain=0:1,samples=100]{0.9*sin(deg(2*pi*x))};
    \addplot[AccentBlue,name path=up,domain=0:1,samples=100]{0.9*sin(deg(2*pi*x))+0.9};
    \addplot[AccentBlue,name path=lo,domain=0:1,samples=100]{0.9*sin(deg(2*pi*x))-0.9};
    \addplot[AccentBlue,fill opacity=0.2,draw=none] fill between[of=up and lo];
    \addplot[only marks,mark=*,mark size=2pt,black]
      coordinates{(0.1,0.5)(0.35,0.65)(0.55,-0.85)(0.85,0.5)};
    \end{axis}
  \end{tikzpicture}
  \end{center}
  \column{.46\textwidth}
  \textbf{With $n=4$ data points:}
  \begin{itemize}
    \item ``S'' shape of sine begins to emerge.
    \item Uncertainty band \alert{narrows} near training points.
    \item Wider in gaps between observations.
    \item Sampled functions becoming more consistent.
  \end{itemize}
  \end{columns}
\end{frame}

\begin{frame}{Bayesian Predictive Distribution — $n=25$}
  \begin{columns}[T]
  \column{.50\textwidth}
  \begin{center}
  \begin{tikzpicture}
    \begin{axis}[
      width=\textwidth,height=5.2cm,
      xlabel={$x$}, ylabel={$t$},
      xmin=0,xmax=1,ymin=-1.8,ymax=1.8,
      grid=both,grid style={line width=.2pt,draw=gray!25},
      tick label style={font=\tiny},thick,
      title={$n=25$: Predictive distribution},
      title style={font=\small}
    ]
    \addplot[MainBlue,domain=0:1,samples=100]{sin(deg(2*pi*x))};
    \addplot[AlertRed,thick,domain=0:1,samples=100]{sin(deg(2*pi*x))};
    \addplot[AccentBlue,name path=up,domain=0:1,samples=100]{sin(deg(2*pi*x))+0.25};
    \addplot[AccentBlue,name path=lo,domain=0:1,samples=100]{sin(deg(2*pi*x))-0.25};
    \addplot[AccentBlue,fill opacity=0.25,draw=none] fill between[of=up and lo];
    \addplot[only marks,mark=*,mark size=1.5pt,black]
      coordinates{(0.04,0.22)(0.08,0.44)(0.12,0.65)(0.16,0.9)(0.20,1.1)
                  (0.24,0.95)(0.28,0.78)(0.32,0.5)(0.36,0.22)(0.40,0.0)
                  (0.44,-0.3)(0.48,-0.62)(0.52,-0.95)(0.56,-0.88)(0.60,-0.77)
                  (0.64,-0.55)(0.68,-0.3)(0.72,0.0)(0.76,0.28)(0.80,0.62)
                  (0.84,0.9)(0.88,1.0)(0.92,0.82)(0.96,0.55)(1.00,0.18)};
    \end{axis}
  \end{tikzpicture}
  \end{center}
  \column{.46\textwidth}
  \textbf{With $n=25$ data points:}
  \begin{itemize}
    \item Predictive mean = excellent approximation of $\sin(2\pi x)$.
    \item Uncertainty band nearly uniform and very narrow.
    \item Sampled curves nearly indistinguishable.
    \item \alert{Posterior has become very concentrated.}
  \end{itemize}
  \vspace{.3em}
  Bayesian approach provides not just a curve, but a ``margin of error'' that naturally shrinks with more data.
  \end{columns}
\end{frame}

\begin{frame}{Predictive Uncertainty Near Training Points}
  \begin{block}{Why uncertainty is localised}
    The precision matrix is updated as:
    \[
      \VSigma_n^{-1} = \VSigma_0^{-1} + \beta\sum_{i=1}^n \ox_i\ox_i^T
    \]
    \begin{itemize}
      \item Near training point $\x_i$: new $\ox$ aligns with well-informed direction $\Rightarrow$ $\ox^T\VSigma\ox$ \alert{small}.
      \item In data-void region: $\ox$ in under-constrained direction $\Rightarrow$ $\ox^T\VSigma\ox$ \alert{large}.
    \end{itemize}
  \end{block}
  \vspace{.3em}
  At $n=25$ with uniform sampling, all basis directions well-covered $\Rightarrow$ nearly uniform uncertainty.
\end{frame}

% =============================================================================
\section{Hyperparameter Marginalisation}
% =============================================================================

\subsection{Fully Bayesian Hyperparameter Treatment}

\begin{frame}{Marginalisation over Hyperparameters}
  \begin{block}{Ideal fully Bayesian integral}
    \[
      p(t|\x,\X,\t) = \int p(t|\x,\ow;\beta)\,p(\ow|\X,\t;\alpha,\beta)\,p(\alpha,\beta|\X,\t)\;d\ow\,d\alpha\,d\beta
    \]
  \end{block}
  \begin{itemize}
    \item Non-linear dependence of $\mathbf{m}_N,\VSigma$ on $\alpha,\beta$ $\Rightarrow$ \alert{analytically intractable}.
    \item Three approximation strategies:
  \end{itemize}
  \vspace{.3em}
  \begin{columns}[T]
  \column{.32\textwidth}
  \begin{block}{Laplace Approx.}
    Gaussian approx.\ around MAP mode via second-order Taylor expansion.
  \end{block}
  \column{.32\textwidth}
  \begin{block}{Variational Inference}
    Cast inference as optimisation; minimise KL divergence between simple proxy $q$ and true posterior.
  \end{block}
  \column{.32\textwidth}
  \begin{block}{MCMC}
    Stochastic sampling; average over samples:
    $p(t|\x)\approx\frac{1}{L}\sum_l p(t|\x,\ow^{(l)},\beta^{(l)})$
  \end{block}
  \end{columns}
\end{frame}

\begin{frame}{Approximation Methods Compared}
  \begin{center}
  \begin{tabular}{lccc}
  \toprule
  \textbf{Method} & \textbf{Type} & \textbf{Cost} & \textbf{Accuracy}\\
  \midrule
  Laplace & Deterministic & Low & Moderate\\
  Variational Inference & Deterministic & Medium & Moderate–High\\
  MCMC & Stochastic & High & High (as $L\to\infty$)\\
  Evidence Framework & Deterministic & Low & Good (practical)\\
  \bottomrule
  \end{tabular}
  \end{center}
  \vspace{.4em}
  \begin{itemize}
    \item In linear regression: \alert{Evidence Framework} is usually sufficient.
    \item For non-linear / safety-critical: MCMC or VI preferred.
    \item MCMC converges to exact integral as $L\to\infty$.
  \end{itemize}
\end{frame}

% =============================================================================
\section{The Evidence Framework (Empirical Bayes)}
% =============================================================================

\subsection{Log-Evidence and Effective Parameters}

\begin{frame}{Evidence Framework — Motivation}
  \begin{alertblock}{Problem}
    Choosing $\alpha$ and $\beta$ by cross-validation is expensive and requires separate validation data.
  \end{alertblock}
  \vspace{.3em}
  \begin{block}{Solution: Empirical Bayes}
    Choose $\alpha,\beta$ to maximise the \alert{marginal likelihood} (evidence):
    \[
      p(\t|\alpha,\beta) = \int p(\t|\ow,\beta)\,p(\ow|\alpha)\,d\ow
    \]
    Integrating out $\ow$ provides an \alert{automatic} trade-off between fit and complexity.
  \end{block}
\end{frame}

\begin{frame}{Log-Evidence Formula}
  Both likelihood and prior are Gaussian $\Rightarrow$ integral is analytic:
  \[
    \log p(\t|\alpha,\beta)
    = \frac{m}{2}\log\alpha + \frac{n}{2}\log\beta - E(\mathbf{m}_p)
    - \frac{1}{2}\log|\mathbf{A}| - \frac{n}{2}\log(2\pi)
  \]
  Key quantities:
  \begin{itemize}
    \item $\mathbf{A} = \VSigma^{-1} = \alpha\I + \beta\oX^T\oX$ — Hessian of error function (posterior curvature).
    \item $E(\mathbf{m}_p) = \tfrac{\beta}{2}\|\t-\oX\mathbf{m}_p\|^2 + \tfrac{\alpha}{2}\mathbf{m}_p^T\mathbf{m}_p$ — regularised error at posterior mean.
  \end{itemize}
  \vspace{.2em}
  \begin{exampleblock}{Automatic complexity control}
    First two terms favour large $\alpha,\beta$; remaining terms penalise complexity.\\
    Maximising evidence $\Leftrightarrow$ automatic bias-variance trade-off.
  \end{exampleblock}
\end{frame}

\begin{frame}{Effective Number of Parameters $\gamma$}
  Let $\eta_i$ be eigenvalues of $\beta\oX^T\oX$. Define:
  \[
    \lambda_i = \frac{\eta_i}{\alpha+\eta_i} \in [0,1]
  \]
  \begin{columns}[T]
  \column{.50\textwidth}
  \begin{block}{Interpretation of $\lambda_i$}
    \begin{itemize}
      \item $\eta_i\gg\alpha$: data dominates $\Rightarrow$ $\lambda_i\to 1$ (parameter data-informed).
      \item $\eta_i\ll\alpha$: prior dominates $\Rightarrow$ $\lambda_i\to 0$ (parameter shrunk to zero).
    \end{itemize}
  \end{block}
  \column{.46\textwidth}
  \begin{block}{Effective parameter count}
    \[
      \gamma = \sum_{i=1}^m \frac{\eta_i}{\alpha+\eta_i}
    \]
    Counts directions in parameter space truly informed by data.
  \end{block}
  \end{columns}
  \vspace{.3em}
  \begin{itemize}
    \item $\gamma\approx m$: data determines nearly all parameters; regularisation minor.
    \item $\gamma\ll m$: most parameters constrained by prior; model uses fewer effective degrees of freedom.
  \end{itemize}
\end{frame}

\begin{frame}{Spring Analogy for $\gamma$}
  \begin{center}
  \begin{tikzpicture}[scale=0.9]
    % Weight node
    \filldraw[MainBlue] (0,0) circle(0.25) node[white,font=\small]{$w_i$};
    % Prior spring (left)
    \draw[thick,decorate,decoration={coil,aspect=0.4,segment length=4pt,amplitude=4pt}]
      (-3.5,0)--(-0.25,0);
    \node[font=\small,above] at (-2,0.4) {Prior spring};
    \node[font=\small,below] at (-2,-0.35) {stiffness $\alpha$};
    \node[font=\small] at (-3.8,0) {$0$};
    % Data spring (right)
    \draw[thick,decorate,decoration={coil,aspect=0.4,segment length=4pt,amplitude=4pt}]
      (0.25,0)--(3.5,0);
    \node[font=\small,above] at (2,0.4) {Data spring};
    \node[font=\small,below] at (2,-0.35) {stiffness $\eta_i$};
    \node[font=\small] at (3.8,0) {$\hat{w}_i$};
    % anchors
    \draw[thick] (-3.5,-0.3)--(-3.5,0.3);
    \draw[thick] (3.5,-0.3)--(3.5,0.3);
  \end{tikzpicture}
  \end{center}
  \vspace{.3em}
  \begin{itemize}
    \item Each weight $w_i$ is pulled by two competing springs.
    \item Prior spring (stiffness $\alpha$) pulls toward \textbf{zero}.
    \item Data spring (stiffness $\eta_i$) pulls toward \textbf{data estimate} $\hat{w}_i$.
    \item $\gamma$ = total number of weights significantly pulled away from zero by data.
  \end{itemize}
\end{frame}

\subsection{Re-estimation of $\alpha$ and $\beta$}

\begin{frame}{Optimal $\alpha$ and $\beta$ — Re-estimation Equations}
  Setting derivatives of log-evidence to zero:
  \vspace{.2em}
  \begin{columns}[T]
  \column{.50\textwidth}
  \begin{block}{Optimal $\alpha$}
    \[
      \alpha = \frac{\gamma}{\mathbf{m}_p^T\mathbf{m}_p}
    \]
    \begin{itemize}
      \item Large posterior weights $\Rightarrow$ reduce $\alpha$ (weaker regularisation).
      \item Small posterior weights $\Rightarrow$ increase $\alpha$ (prevent overfitting).
    \end{itemize}
  \end{block}
  \column{.46\textwidth}
  \begin{block}{Optimal $\beta$}
    \[
      \frac{1}{\beta} = \frac{1}{n-\gamma}\sum_{i=1}^n(t_i-\mathbf{m}_p^T\ox_i)^2
    \]
    Residual variance corrected by effective DOF.\\
    $n-\gamma$ = effective sample size.
  \end{block}
  \end{columns}
\end{frame}

\begin{frame}{Iterative Re-estimation Algorithm}
  Since $\mathbf{m}_p$ and $\gamma$ both depend on $\alpha$ and $\beta$, solve iteratively:
  \vspace{.3em}
  \begin{enumerate}
    \item \textbf{Initialize} $\alpha$ and $\beta$ (e.g.\ $\alpha=1, \beta=1$).
    \item \textbf{Compute} posterior mean $\mathbf{m}_p$ and covariance $\VSigma_p$.
    \item \textbf{Compute} effective parameters $\gamma = \sum_i \tfrac{\eta_i}{\alpha+\eta_i}$.
    \item \textbf{Update} $\alpha = \gamma/(\mathbf{m}_p^T\mathbf{m}_p)$ and $\beta^{-1} = \tfrac{1}{n-\gamma}\sum_i(t_i-\mathbf{m}_p^T\ox_i)^2$.
    \item \textbf{Repeat} steps 2–4 until convergence.
  \end{enumerate}
  \vspace{.4em}
  \begin{exampleblock}{Advantages}
    \begin{itemize}
      \item Converges quickly in practice.
      \item Relies solely on training data — \alert{no validation set needed}.
      \item Automatically adapts to data complexity.
    \end{itemize}
  \end{exampleblock}
\end{frame}

\begin{frame}{Evidence Framework — Summary}
  \begin{block}{What it achieves}
    \begin{itemize}
      \item Noisy data ($\small\beta$): evidence maximised by large $\alpha$ $\Rightarrow$ weights shrunk toward zero.
      \item Clean, informative data: $\gamma\to m$ $\Rightarrow$ model uses full expressive power.
      \item Self-consistent: no external tuning or cross-validation required.
    \end{itemize}
  \end{block}
  \vspace{.4em}
  \begin{center}
  \begin{tikzpicture}[
    every node/.style={font=\small},
    box/.style={draw=MainBlue,fill=torgray,rounded corners=4pt,
                minimum width=3.2cm,minimum height=1cm,align=center}
  ]
    \node[box] (init) {Initialise\\$\alpha,\beta$};
    \node[box,right=1cm of init] (post) {Compute\\$\mathbf{m}_p,\VSigma_p$};
    \node[box,right=1cm of post] (gam) {Compute $\gamma$};
    \node[box,below=0.6cm of gam] (upd) {Update\\$\alpha,\beta$};
    \node[box,left=1cm of upd] (conv) {Converged?};
    \node[box,left=1cm of conv,fill=ExGreen!20,draw=ExGreen] (done) {Done:\\output $\alpha^*,\beta^*$};
    \draw[->,thick,MainBlue] (init)--(post);
    \draw[->,thick,MainBlue] (post)--(gam);
    \draw[->,thick,MainBlue] (gam)--(upd);
    \draw[->,thick,MainBlue] (upd)--(conv);
    \draw[->,thick,MainBlue] (conv)--node[above,font=\tiny]{no} (post);
    \draw[->,thick,ExGreen] (conv)--node[above,font=\tiny]{yes} (done);
  \end{tikzpicture}
  \end{center}
\end{frame}

% =============================================================================
\section{Summary}
% =============================================================================

\begin{frame}{Chapter Summary — Linear Models}
  \begin{columns}[T]
  \column{.50\textwidth}
  \begin{block}{Core concepts}
    \begin{itemize}
      \item Linear predictor: $h(\x;\w,b)=\w^T\x+b$.
      \item Basis functions expand expressive power without breaking parameter linearity.
      \item Design matrix $\VPhi$: transforms dataset.
    \end{itemize}
  \end{block}
  \begin{block}{Optimisation}
    \begin{itemize}
      \item Normal equations: closed-form, $O(d^3)$.
      \item Gradient descent: iterative, scalable.
    \end{itemize}
  \end{block}
  \column{.46\textwidth}
  \begin{block}{Regularisation}
    \begin{itemize}
      \item Ridge ($q{=}2$): shrinks weights, closed-form.
      \item Lasso ($q{=}1$): sparse solution, feature selection.
      \item $\lambda$ controls bias-variance trade-off.
    \end{itemize}
  \end{block}
  \begin{block}{Probabilistic / Bayesian}
    \begin{itemize}
      \item MLE $\equiv$ Least Squares (Gaussian noise).
      \item MAP $\equiv$ Ridge ($\lambda{=}\alpha/\beta$).
      \item Full Bayes: predictive distribution with aleatoric + epistemic uncertainty.
      \item Evidence Framework: automatic hyperparameter tuning.
    \end{itemize}
  \end{block}
  \end{columns}
\end{frame}

\begin{frame}{Key Equations at a Glance}
  \begin{center}
  \renewcommand{\arraystretch}{1.6}
  \begin{tabular}{ll}
  \toprule
  \textbf{Concept} & \textbf{Formula}\\
  \midrule
  Normal equations & $\ow^* = (\oX^T\oX)^{-1}\oX^T\t$\\
  Ridge solution & $\ow = (\lambda\I+\oX^T\oX)^{-1}\oX^T\t$\\
  Posterior mean & $\mathbf{m}_p = \beta\VSigma_p\oX^T\t$\\
  Posterior cov. & $\VSigma_p = (\alpha\I+\beta\oX^T\oX)^{-1}$\\
  Predictive mean & $m(\x)=\beta\ox^T\VSigma\oX^T\t$\\
  Predictive var. & $\sigma^2(\x)=\beta^{-1}+\ox^T\VSigma\ox$\\
  Effective params & $\gamma=\sum_i\tfrac{\eta_i}{\alpha+\eta_i}$\\
  Optimal $\alpha$ & $\alpha=\gamma/(\mathbf{m}_p^T\mathbf{m}_p)$\\
  \bottomrule
  \end{tabular}
  \end{center}
\end{frame}

%\end{document}
